\section{Redovisningen du behöver}
\label{sec:redovisning}

Värdering börjar i bokslutet. Du behöver inte kunna bokföra, men du måste kunna läsa de tre
rapporterna och veta hur en ändring i en av dem slår igenom i de andra två. Talen i exemplen är
Atlas Copcos för 2025 så som datalagret levererar dem, i mdkr.

\subsection{De tre rapporterna}

\begin{description}
  \item[Resultaträkningen] \eng{income statement} visar intäkter och kostnader under ett år.
    Den är periodiserad: en kostnad tas när resursen förbrukas, inte när den betalas.
  \item[Balansräkningen] \eng{balance sheet} visar vad bolaget äger och hur det är finansierat
    en viss dag. Den ska alltid gå ihop:
    $\text{tillgångar} = \text{skulder} + \text{eget kapital}$.
  \item[Kassaflödesanalysen] \eng{cash flow statement} visar vart pengarna faktiskt tog vägen
    under året, uppdelat på löpande verksamhet, investeringar och finansiering.
\end{description}

\begin{table}[ht]
\centering\small
\caption{Resultaträkningens trappa, Atlas Copco 2025 (mdkr). Källa: projektets datalager.}
\label{tab:rr}
\begin{tabular}{@{}llr@{}}
\toprule
Rad & Engelska & 2025 \\
\midrule
Omsättning & \en{revenue}, \en{sales} & 168,3 \\
$-$ Kostnad för sålda varor & \en{cost of goods sold, COGS} & 96,1 \\
$=$ Bruttoresultat & \en{gross profit} & 72,2 \\
$-$ Övriga rörelsekostnader (före avskrivningar) & \en{operating expenses} & 28,1 \\
$=$ \textbf{EBITDA} & & 44,1 \\
$-$ Avskrivningar & \en{depreciation and amortization, D\&A} & 9,5 \\
$=$ \textbf{EBIT}, rörelseresultat & \en{operating income} & 34,6 \\
$-$ Räntekostnader & \en{interest expense} & 0,9 \\
$=$ Resultat före skatt & \en{pretax income, EBT} & 33,7 \\
$-$ Skatt & \en{income tax} & 7,2 \\
$=$ Nettoresultat & \en{net income} & 26,4 \\
\bottomrule
\end{tabular}
\end{table}

Balansräkningen samma år: tillgångar 202,5, skulder 92,1, eget kapital hänförligt till
aktieägarna 110,2 och minoritetens andel 0,2. Kassaflödet från den löpande verksamheten var
32,6, investeringarna i anläggningar (capex) 6,2 och utdelningen 14,6.

\paragraph{Hur de hänger ihop.} Figur~\ref{fig:tre-rapporter} visar kopplingarna. Lär dig dem som
en kedja:

\begin{enumerate}
  \item Nettoresultatet är första raden i kassaflödesanalysen och ökar det egna kapitalet
    (balanserade vinstmedel, \en{retained earnings}), minus utdelning.
  \item Avskrivningar är en kostnad utan utbetalning. De läggs tillbaka i kassaflödet och minskar
    anläggningstillgångarna i balansräkningen.
  \item Investeringar (capex) syns inte i resultaträkningen. De är ett utflöde i kassaflödet och
    ökar anläggningstillgångarna.
  \item Ökat rörelsekapital binder pengar: ett utflöde i kassaflödet, en ökning av
    omsättningstillgångarna i balansräkningen.
  \item Summan av de tre kassaflödena är årets förändring av kassan, som är en rad i
    balansräkningen. Där sluts cirkeln.
\end{enumerate}

\begin{figure}[ht]
\centering
\begin{tikzpicture}[font=\small\sffamily,
  box/.style={draw=accent,rounded corners=2pt,align=left,inner sep=5pt,text width=3.9cm,
              fill=accent!4},
  head/.style={font=\small\sffamily\bfseries,text=accent},
  link/.style={-{Stealth[length=5pt]},thick,draw=black!60},
  lbl/.style={font=\footnotesize\sffamily,fill=white,inner sep=1.5pt,text=black!75}]
  \node[box] (rr) {\textbf{Resultaträkning}\\ Omsättning\\ EBITDA\\ $-$ Avskrivningar\\ EBIT\\
    $-$ Ränta, skatt\\ \textbf{Nettoresultat}};
  \node[box,right=1.5cm of rr] (kf) {\textbf{Kassaflödesanalys}\\ Nettoresultat\\
    $+$ Avskrivningar\\ $-$ Ökning rörelsekapital\\ $-$ Capex\\ $\pm$ Lån, utdelning\\
    \textbf{Förändring av kassa}};
  \node[box,right=1.5cm of kf] (br) {\textbf{Balansräkning}\\ Kassa\\ Rörelsekapital\\
    Anläggningstillgångar\\ Skulder\\ Eget kapital\\ \emph{(tillgångar = skulder + EK)}};
  \draw[link] (rr.east) to node[lbl,above=1pt] {resultat} (kf.west);
  \draw[link] (kf.east) to node[lbl,above=1pt] {kassa} (br.west);
  \draw[link] (rr.south) to[out=-90,in=-90,looseness=0.42]
    node[lbl,pos=0.5] {nettoresultat ökar eget kapital, avskrivningar minskar anläggningstillgångar}
    (br.south);
\end{tikzpicture}
\caption{De tre rapporterna och hur de hänger ihop.}
\label{fig:tre-rapporter}
\end{figure}

\subsection{EBITDA}

EBITDA \eng{earnings before interest, taxes, depreciation and amortization} är rörelseresultatet
före avskrivningar. Måttet är populärt av tre skäl. Det påverkas inte av hur bolaget är
finansierat (räntan ligger längre ned), inte av avskrivningsprinciper, och det är en grov
approximation av kassaflödet från rörelsen före investeringar.

\begin{vanligtfel}
Att behandla EBITDA som kassaflöde. EBITDA tar inte hänsyn till investeringar, bundet
rörelsekapital eller skatt. Två bolag med samma EBITDA kan ha helt olika fritt kassaflöde.
Sedan redovisningsregeln IFRS~16 ligger dessutom hyreskostnader för leasade tillgångar som
avskrivning och ränta, alltså \emph{under} EBITDA. Motorn justerar inte för det, och
kassaflödena under prognosåren blir därför något för höga för bolag med mycket leasing.
\end{vanligtfel}

\subsection{Rörelsekapital}

Rörelsekapital \eng{net working capital, NWC} är de pengar som är bundna i den dagliga
verksamheten. Motorn använder den smala, operativa definitionen:
\begin{equation}
  \NWC = \text{kundfordringar} + \text{lager} - \text{leverantörsskulder}.
  \label{eq:nwc}
\end{equation}
För Atlas Copco 2025: $29,9 + 26,7 - 16,4 = 40,2$~mdkr, eller 23,9\pct{} av omsättningen. När
bolaget växer måste rörelsekapitalet växa med, och den ökningen är ett kassautflöde. I prognosen
antar motorn att rörelsekapitalet är en fast andel $\nu$ av omsättningen, så att
$\Delta\NWC_t = \nu\,(R_t - R_{t-1})$. Ett bolag som får betalt i förskott kan ha negativt
rörelsekapital, och då \emph{frigör} tillväxt pengar.

\subsection{Fritt kassaflöde: FCFF och FCFE}

Fritt kassaflöde är det som blir över när rörelsen har fått sitt. Det finns två varianter, och
skillnaden är vem pengarna tillhör.

\begin{description}
  \item[FCFF] \eng{free cash flow to the firm, unlevered free cash flow} är kassaflödet till
    \emph{alla} finansiärer, långivare och aktieägare, före räntor:
    \begin{equation}
      \FCFF = {\underbrace{\EBIT\,(1-\tau)}_{\NOPAT}} + \text{avskrivningar} - \text{capex}
              - \Delta\NWC .
      \label{eq:fcff}
    \end{equation}
    Skatten räknas som om bolaget saknade skuld. Värdet av ränteavdraget kommer i stället in via
    diskonteringsräntan (avsnitt~\ref{sec:wacc}).
  \item[FCFE] \eng{free cash flow to equity, levered free cash flow} är det som blir kvar till
    aktieägarna sedan långivarna fått sitt:
    \begin{equation}
      \FCFE = \FCFF - \text{ränta}\,(1-\tau) + \text{nettoupplåning}.
      \label{eq:fcfe}
    \end{equation}
\end{description}

Regeln är att kassaflöde och diskonteringsränta måste höra ihop. FCFF diskonteras med WACC och
ger rörelsevärdet. FCFE diskonteras med avkastningskravet på eget kapital och ger aktievärdet
direkt. Motorn räknar bara med FCFF, vilket också är standard på en bank, eftersom man då inte
behöver prognostisera skuldsättningen år för år.

\paragraph{Exempel.} Motorns första prognosår för Atlas Copco (2026): EBITDA 47,10 och
avskrivningar 9,00 ger EBIT 38,11. Skatt 20,6\pct{} på det är 7,85, så $\NOPAT = 30,26$. Capex är
6,09 och rörelsekapitalet ökar med 2,54. Enligt~\eqref{eq:fcff}:
$\FCFF = 30,26 + 9,00 - 6,09 - 2,54 \approx 30,6$~mdkr.

\subsection{Klassikern: avskrivningarna ökar med 10 kronor}
\label{sec:avskrivning10}

Frågan kommer i nästan varje intervju, och den testar om du kan kedjan ovan. Gå alltid i
ordningen resultaträkning, kassaflöde, balansräkning, och avsluta med att visa att
balansräkningen går ihop. Med projektets skattesats 20,6\pct:

\begin{table}[H]
\centering\small
\caption{Avskrivningarna ökar med 10\kr. Skattesats 20,6\pct.}
\label{tab:avskrivning}
\begin{tabular}{@{}llr@{}}
\toprule
Rapport & Rad & Förändring \\
\midrule
Resultaträkning & Resultat före skatt & $-10,00$ \\
                & Skatt ($20,6\,\%$ av $-10$) & $+2,06$ \\
                & Nettoresultat & $-7,94$ \\
\addlinespace
Kassaflödesanalys & Nettoresultat & $-7,94$ \\
                  & Återlagd avskrivning (ingen utbetalning) & $+10,00$ \\
                  & Kassaflöde från rörelsen, och därmed kassan & $+2,06$ \\
\addlinespace
Balansräkning & Anläggningstillgångar & $-10,00$ \\
              & Kassa & $+2,06$ \\
              & Summa tillgångar & $-7,94$ \\
              & Eget kapital (balanserat resultat) & $-7,94$ \\
\bottomrule
\end{tabular}
\end{table}

Kassan \emph{ökar} alltså, fast vinsten minskar. Avskrivningen är ingen utbetalning, men den
sänker skatten med $0,206 \cdot 10 = 2,06$\kr. Det kallas avskrivningarnas skattesköld
\eng{depreciation tax shield}.

\begin{intervjutips}
Säg skattesatsen högt innan du börjar räkna (”I'll assume a 20.6 percent tax rate, the Swedish
corporate rate”) och avsluta med ”and the balance sheet balances”. Amerikanska läroböcker räknar
ofta med 40 eller 25\pct. Metoden är densamma, och med 20\pct{} blir huvudräkningen enkel:
$-8$, $+2$, $-8$.
\end{intervjutips}
